\documentclass[11pt,a4paper]{article}
\usepackage[utf8]{inputenc}
\usepackage[T1]{fontenc}
\usepackage{amsmath,amssymb}
\usepackage{booktabs}
\usepackage{multirow}
\usepackage{graphicx}
\usepackage{xcolor}
\usepackage{hyperref}
\usepackage[margin=2.5cm]{geometry}
\usepackage{natbib}
\usepackage{enumitem}
\usepackage{array}

\newcommand{\fm}[1]{\textbf{FM#1}}
\newcommand{\pp}{\,pp}
\newcommand{\ci}[2]{[#1,\;#2]}

\title{Why Calibration Fails: A Taxonomy of Failure Modes\\in Dialect-Aware LLM User Simulation}

\author{Anonymous}

\date{}

\begin{document}
\maketitle

% ============================================================
\begin{abstract}
LLM-based user simulators exhibit systematic dialect-dependent behavioral biases---speakers of African American Vernacular English (AAVE) and Indian English receive measurably different politeness and question-asking patterns compared to Standard American English.
We attempt to correct these biases through two complementary calibration strategies: Behavioral Calibration Prompting (BCP) at inference time and Direct Preference Optimization (DPO) at training time.
Both fail.
Through controlled experiments on Qwen2.5-7B-Instruct and Qwen2.5-14B-Instruct in the $\tau$-bench retail domain, we identify four distinct failure mechanisms: (1)~global behavior suppression, where calibration prompts collapse task completion; (2)~cross-dimension interference, where targeting one behavioral dimension perturbs others; (3)~overcorrection rebound, where extreme calibration targets trigger opposite-direction shifts; and (4)~DPO--BCP interaction reversal, where preference-optimized models reverse the direction of prompt-based calibration.
All four mechanisms are more severe at 7B than 14B, with bootstrap confidence intervals confirming statistical significance.
Yet calibration remains ineffective at both scales---the H2 gap-reduction criterion fails across all conditions.
Scaling changes \emph{how} calibration fails, not \emph{whether} it fails.
These findings establish concrete constraints for future calibration research and reveal a safety--efficacy gap: scaling eliminates destructive side effects without improving calibration effectiveness.
\end{abstract}

% ============================================================
\section{Introduction}
\label{sec:intro}

LLM-based user simulators are increasingly used to evaluate task-oriented dialogue agents \citep{yao2024taubench}, yet recent work reveals a troubling sim-to-real gap: simulated users exhibit systematic behavioral biases that vary with the dialect of the user persona \citep{sim2real2026gap, seshadri2026lost}.
AAVE and Indian English speakers receive different politeness levels and question-asking rates than Standard American English speakers, and these biases propagate into downstream agent evaluation metrics.

A natural response is calibration: adjust the simulator's behavioral rates to match empirical targets.
We investigate two complementary paradigms.
\textbf{Behavioral Calibration Prompting (BCP)} injects target behavioral rates directly into the system prompt at inference time.
\textbf{Direct Preference Optimization (DPO)} \citep{rafailov2023dpo} fine-tunes the model on dialect-calibrated preference pairs at training time.
Both are intuitive and represent the two main intervention points---prompt-time and training-time.

Both fail systematically.
Rather than treating this as a negative result, we conduct a controlled empirical analysis to understand \emph{why} calibration fails.
Our main contribution is the first taxonomy of calibration failure modes for LLM user simulators, comprising four mechanisms (\fm{1}--\fm{4}).
We evaluate these mechanisms across two model scales (7B and 14B parameters) within the Qwen2.5 family and find that failure severity is \textbf{scale-dependent}: every mechanism is more severe at 7B than 14B, yet calibration remains ineffective at both scales.

This reveals a \textbf{safety--efficacy gap}: scaling from 7B to 14B eliminates BCP's destructive side effects (completion collapse, overcorrection rebound) but does not improve its calibration effectiveness.
There is no ``sweet spot'' scale where BCP works---smaller models are harmed by it, larger models ignore it.

Our contributions are:
\begin{enumerate}[nosep]
    \item A four-mechanism taxonomy of calibration failure modes (\fm{1}--\fm{4}), each statistically validated with bootstrap confidence intervals.
    \item Evidence that \fm{4} (DPO--BCP interaction reversal) is scale-dependent: strong at 7B (10/12 directional consistency, $p=0.006$) but attenuated to chance at 14B (4/12, CI includes zero), with the attenuation itself statistically significant (14.7\pp, 95\% CI $\ci{8.1}{21.6}$).
    \item A scale-dependent failure severity analysis showing all four mechanisms exhibit capacity-dependent severity, establishing concrete constraints for future calibration research.
\end{enumerate}

% ============================================================
\section{Related Work}
\label{sec:related}

\paragraph{Persona prompting and capability trade-offs.}
\citet{basil2025playing_pretend} show that persona prompting degrades LLM accuracy (MMLU drops from 71.6\% to 66.3\%), and \citet{hu2026expert_persona} document accuracy degradation up to 30\pp{} under persona, with instruction-tuned models being more susceptible.
Our \fm{1} (global behavior suppression) is a task-oriented instantiation of this persona--accuracy trade-off, but we additionally demonstrate its scale dependence.

\paragraph{Bias spillover in alignment.}
\citet{gupta2026intra_fairness} document bias spillover effects in targeted alignment, where correcting one fairness dimension perturbs others.
Our \fm{2} (cross-dimension interference) provides a parallel finding in behavioral calibration.

\paragraph{DPO alignment fragility.}
\citet{lee2024dpo_toxicity} show that DPO bypasses rather than removes capabilities, producing fragile alignment that can be reversed.
Our \fm{4} (DPO--BCP interaction reversal) demonstrates this fragility in a calibration context: DPO training is effective in isolation but reverses calibration direction when combined with BCP prompts.

\paragraph{LLM user simulation.}
$\tau$-bench \citep{yao2024taubench} provides a standardized evaluation framework for tool-calling dialogue agents.
\citet{sim2real2026gap} quantify the sim-to-real gap in $\tau$-bench, and \citet{seshadri2026lost} identify dialect-dependent calibration as a core source of this gap.
We build directly on this evaluation infrastructure.

% ============================================================
\section{Method}
\label{sec:method}

\subsection{Task Setting}

We evaluate in the $\tau$-bench retail domain, a 16-tool customer service environment where a user simulator interacts with a tool-calling agent across multi-turn dialogues.
Each episode involves a user persona with a dialect variant (Standard American English, AAVE, Indian English, or Overcorrected) interacting for task completion.
We measure three behavioral dimensions: \textbf{task completion rate}, \textbf{politeness rate} (fraction of turns containing politeness markers), and \textbf{question rate} (questions per turn).

\subsection{Evaluation Criteria}

We define three hypotheses for successful calibration:
\begin{itemize}[nosep]
    \item \textbf{H1} (Dialect sensitivity): The model produces measurably different behavioral rates across dialect variants ($\geq$5\% relative difference).
    \item \textbf{H2} (Gap reduction): BCP reduces the behavioral gap between dialect variants and the baseline by $\geq$30\%.
    \item \textbf{H3} (No harm): Calibrated variants retain $\geq$80\% of baseline task completion.
\end{itemize}

\subsection{Calibration Strategies}

\paragraph{Behavioral Calibration Prompting (BCP).}
We inject explicit target behavioral rates into the system prompt.
For each dialect variant, BCP specifies the desired politeness rate and question rate, derived from empirical dialect profiles.
Variants include single-dimension targets (\texttt{single\_dim\_politeness}, \texttt{single\_dim\_question}) and multi-dimension targets (AAVE, Indian English, Overcorrected).

\paragraph{Direct Preference Optimization (DPO).}
We fine-tune on 1,434 dialect-calibrated preference pairs, where preferred responses exhibit appropriate dialect-matched behavioral rates.
Training uses LoRA ($r{=}16$) on \texttt{q\_proj}/\texttt{v\_proj} with FP16 precision.
At 7B: single RTX 4090, 3 epochs, $\sim$1.4h.
At 14B: 2$\times$ RTX 4090 with tensor parallelism, 3 epochs, $\sim$7.5h.

\subsection{Statistical Methodology}

All core comparisons use non-parametric bootstrap with $B{=}10{,}000$ episode-level resamples and BCa correction for confidence intervals.
Directional consistency is assessed via sign tests under a binomial null ($p{=}0.5$).

% ============================================================
\section{Experiments and Results}
\label{sec:experiments}

We present results organized by the four failure mechanisms of our taxonomy, followed by cross-scale analysis and combined DPO--BCP evaluation at 14B.

\subsection{Failure Mode Taxonomy}

\subsubsection{\fm{1}: Global Behavior Suppression}
\label{sec:fm1}

BCP calibration prompts at 7B collapse task completion regardless of the target dimension.
Mean completion retention across calibrated variants is 52.3\% (threshold $\geq$80\%), with AAVE completion dropping to 10.0\% (from baseline 36.8\%, a 72.8\% decrease) and Indian English to 15.0\% (59.2\% decrease).
The 7B-vs-14B completion difference is $+24.8\pp$ (95\% CI $\ci{16.4}{33.2}$, excludes zero), confirming this is capacity-dependent.

At 14B, \fm{1} disappears entirely: H3 passes 4/4 runs with all variants retaining $>$97\% of baseline completion.
The 14B model maintains a mean completion of 61.6\% versus 36.8\% at 7B, a 67\% improvement from scaling alone.

\fm{1} parallels the persona--accuracy trade-off documented by \citet{basil2025playing_pretend} and \citet{hu2026expert_persona}: BCP prompts function as behavioral personas that overwhelm the 7B model's limited capacity, degrading core task performance.

\subsubsection{\fm{2}: Cross-Dimension Interference}
\label{sec:fm2}

Targeting one behavioral dimension perturbs non-targeted dimensions.
When targeting question rate (\texttt{single\_dim\_question}), politeness drops by $-25.6\%$ relative to baseline (95\% CI $\ci{-30.1}{-20.8}$, excludes zero; $n{=}20$).
The absolute shift is $-14.3\pp$ (from 0.559 to 0.416).

In the reverse direction, targeting politeness (\texttt{single\_dim\_politeness}) reduces question rate by 52.3\% relative.
We note that the baseline question rate is low (0.044, approximately 0.8 questions per task), so the 52.3\% relative drop corresponds to an absolute change from 0.044 to 0.021 ($-2.3\pp$).
We report this arm descriptively with absolute values.

At 14B, \fm{2} persists but is attenuated: politeness interference drops to $-8.0\%$ relative (from $-25.6\%$ at 7B).
At 14B with DPO, the interference is $-3.8\pp$ in politeness and $-1.9\pp$ in question rate---weakened but nonzero.
\fm{2} is the only mechanism that persists across both scales and both calibration paradigms (BCP and DPO).

\subsubsection{\fm{3}: Overcorrection Rebound}
\label{sec:fm3}

Requesting an extreme behavioral shift (the Overcorrected variant) causes the model to rebound in the opposite direction at 7B.
The overcorrected variant produces politeness $+7.23\pp$ \emph{above} baseline (95\% CI $\ci{3.8}{10.6}$, excludes zero), despite being targeted to dramatically reduce politeness.
This effect is consistent across all 4 runs (4/4 positive direction), with per-run rebounds ranging from $+4.2\pp$ to $+10.4\pp$.

At 14B, \fm{3} disappears: 0/4 runs show rebound, and the overcorrected variant's politeness is consistently \emph{below} baseline.
This suggests that 7B models have an ``attractor basin'' for default behavioral rates that resists large perturbations, while 14B models have sufficient capacity to either obey or simply ignore the calibration instruction without rebounding.

\subsubsection{\fm{4}: DPO--BCP Interaction Reversal}
\label{sec:fm4}

DPO training at 7B is effective in isolation: rewards accuracy reaches 0.75--1.0, loss converges to 0.653, and task completion improves from 40\% to 60\% after weight merging.
However, when combined with BCP prompts, DPO \emph{reverses} the calibration direction.

Under base BCP (without DPO), the AAVE politeness gap is $-9.1\pp$ (correct direction---calibration reduces the gap).
Under DPO+BCP, this gap flips to $+6.4\pp$ (\textbf{reversed}---calibration now \emph{increases} the gap).
The reversal swing is $+15.8\pp$ (95\% CI $\ci{10.2}{21.8}$, $p<0.01$).
Table~\ref{tab:fm4_direction} shows that 10/12 variant$\times$run cells exhibit gap expansion, with a sign test $p{=}0.006$ under the binomial null.

\begin{table}[t]
\centering
\caption{FM4 per-cell gap direction at 7B. ``$\uparrow$ expanded'' indicates DPO+BCP \emph{worsened} the politeness gap relative to base BCP. Gap change percentages are relative to base BCP gap magnitude.}
\label{tab:fm4_direction}
\small
\begin{tabular}{lccccr}
\toprule
\textbf{Variant} & \textbf{Run 1} & \textbf{Run 2} & \textbf{Run 3} & \textbf{Run 4} & \textbf{Row $p$} \\
\midrule
AAVE & $\uparrow$ ($-48.7\%$) & $\uparrow$ ($-35.1\%$) & $\uparrow$ ($-15.2\%$) & $\downarrow$ ($+5.9\%$) & 0.023 \\
Indian & $\uparrow$ ($-37.3\%$) & $\uparrow$ ($-25.9\%$) & $\downarrow$ ($+6.8\%$) & $\uparrow$ ($-31.1\%$) & 0.023 \\
Overcorrected & $\uparrow$ ($-25.2\%$) & $\uparrow$ ($-11.6\%$) & $\uparrow$ ($-11.3\%$) & $\uparrow$ ($-23.9\%$) & 0.002 \\
\midrule
Column total & 3/3 $\uparrow$ & 3/3 $\uparrow$ & 2/3 $\uparrow$ & 2/3 $\uparrow$ & \textbf{10/12} \\
\bottomrule
\end{tabular}
\end{table}

This reversal is not a DPO training failure---DPO genuinely improves tool-calling ability.
Rather, it is an \textbf{interaction effect}: DPO alignment and BCP prompting interact adversarially, consistent with findings that DPO bypasses rather than removes capabilities \citep{lee2024dpo_toxicity}.

% ============================================================
\subsection{Bootstrap Confidence Intervals}
\label{sec:bootstrap}

Table~\ref{tab:bootstrap_ci} summarizes bootstrap confidence intervals for all core effects.
All five comparisons exclude zero, upgrading the taxonomy from directional consistency arguments to statistically confirmed mechanisms.

\begin{table}[t]
\centering
\caption{Bootstrap 95\% BCa confidence intervals for core effects ($B{=}10{,}000$, episode-level resampling).}
\label{tab:bootstrap_ci}
\small
\begin{tabular}{lcccc}
\toprule
\textbf{Comparison} & \textbf{Observed} & \textbf{95\% BCa CI} & \textbf{Excl.\ zero?} & \textbf{$n$} \\
\midrule
FM4: Gap reversal swing (pol \pp) & $+15.8\pp$ & $\ci{+10.2}{+21.8}$ & Yes ($p<0.01$) & 80 vs 19 \\
FM3: Overcorrection rebound (\pp) & $+7.23\pp$ & $\ci{+3.8}{+10.6}$ & Yes & 80 \\
7B-vs-14B completion diff (\pp) & $+24.8\pp$ & $\ci{+16.4}{+33.2}$ & Yes & 80 vs 19 \\
FM2: Pol non-target shift (\% rel) & $-25.6\%$ & $\ci{-30.1}{-20.8}$ & Yes & 20 \\
FM4: Directional consistency & 10/12 & med.\ 10, $\ci{8}{12}$ & $p{=}0.006$ & 12 cells \\
\bottomrule
\end{tabular}
\end{table}

The \fm{4} reversal has the strongest signal (CI lower bound at $+10.2\pp$, more than $5\times$ the margin from zero).
\fm{3} rebound is weakest in absolute terms but still significant (lower bound $+3.8\pp$).
The scale completion difference is unambiguous ($+24.8\pp$, lower bound $+16.4\pp$).
The \fm{2} politeness interference CI is tight ($\ci{-30.1}{-20.8}$) despite $n{=}20$ because the effect size is large ($>$2 SD shift with minimal overlap between conditions).

% ============================================================
\subsection{Scale-Dependent FM4: 14B DPO+BCP}
\label{sec:14b_dpo}

To test whether \fm{4} is structural or capacity-dependent, we replicate the DPO+BCP experiment at 14B.
Phase~A: DPO fine-tuning of Qwen2.5-14B-Instruct on the same 1,434 preference pairs (LoRA $r{=}16$, FP16, tensor parallelism across 2$\times$RTX~4090, 3 epochs, $\sim$7.5h).
Phase~B: BCP evaluation of the weight-merged model (4 runs $\times$ 6 variants $\times$ 20 tasks = 480 episodes).

Table~\ref{tab:14b_training} compares training dynamics across scales.
The 14B model achieves comparable loss (0.658 vs 0.653) but slightly lower rewards accuracy (0.70--0.88 vs 0.75--1.0), because the 14B baseline is already closer to the DPO target distribution.

\begin{table}[t]
\centering
\caption{DPO training metrics: 14B vs 7B.}
\label{tab:14b_training}
\small
\begin{tabular}{lcc}
\toprule
\textbf{Metric} & \textbf{14B} & \textbf{7B} \\
\midrule
Final loss & 0.658 & 0.653 \\
Rewards accuracy & 0.70--0.88 & 0.75--1.0 \\
Rewards margin & 0.03--0.12 & 0.05--0.17 \\
Tokenizer mismatch & 0\% & 0\% \\
Training time & $\sim$7.5h (2$\times$4090) & $\sim$1.4h (1$\times$4090) \\
Epochs / Steps & 3 / 1173 & 3 / 1173 \\
\bottomrule
\end{tabular}
\end{table}

Table~\ref{tab:14b_eval} shows per-run evaluation results.
The 14B DPO model passes H1 (dialect sensitivity) and H3 (no harm) in all 4 runs, but fails H2 (gap reduction) in all 4 runs---identical to the pattern at 7B.

\begin{table}[t]
\centering
\caption{14B DPO+BCP per-run evaluation. Politeness gaps are in percentage points relative to baseline.}
\label{tab:14b_eval}
\small
\begin{tabular}{lccccc}
\toprule
 & \textbf{Run 1} & \textbf{Run 2} & \textbf{Run 3} & \textbf{Run 4} & \textbf{Mean} \\
\midrule
Completion & 0.632 & 0.547 & 0.579 & 0.611 & 0.592 \\
Baseline pol & 0.511 & 0.498 & 0.527 & 0.536 & 0.518 \\
AAVE gap (\pp) & $+6.2$ & $-2.1$ & $-0.6$ & $+3.7$ & $+1.8$ \\
Indian gap (\pp) & $+7.5$ & $-0.8$ & $-1.5$ & $+4.4$ & $+2.4$ \\
Overcorr gap (\pp) & $-0.3$ & $-1.2$ & $-0.8$ & $-1.8$ & $-1.0$ \\
H1 / H2 / H3 & P/F/P & P/F/P & P/F/P & P/F/P & 4--0--4 \\
\bottomrule
\end{tabular}
\end{table}

\paragraph{FM4 is scale-dependent.}
Table~\ref{tab:cross_scale_fm4} presents the cross-scale comparison.
At 7B, the pooled reversal is $+15.8\pp$ with 10/12 directional consistency and a CI that excludes zero ($\ci{10.2}{21.8}$).
At 14B, the pooled reversal is only $+1.1\pp$ with 4/12 directional consistency (near chance) and a CI that includes zero ($\ci{-1.6}{5.9}$).
The attenuation itself---14.7\pp---is statistically significant (95\% CI $\ci{8.1}{21.6}$, excludes zero), confirming that FM4 is capacity-dependent.

\begin{table}[t]
\centering
\caption{Cross-scale FM4 reversal comparison. The 14B CI includes zero; the attenuation CI excludes zero.}
\label{tab:cross_scale_fm4}
\small
\begin{tabular}{lccc}
\toprule
\textbf{Metric} & \textbf{7B} & \textbf{14B} & \textbf{Attenuation} \\
\midrule
Pooled reversal (\pp) & $+15.8$ & $+1.1$ & $-14.7\pp$ (93\% reduction) \\
AAVE gap (\pp) & $+6.4$ & $+1.8$ & $-4.6\pp$ \\
Indian gap (\pp) & $+6.3$ & $+2.4$ & $-3.9\pp$ \\
H2 gap reduction & $-10.3\%$ & $-3.1\%$ & $+7.2\pp$ less reversal \\
Directional consistency & 10/12 (83\%) & 4/12 (33\%) & $-50\pp$ drop \\
95\% CI (\pp) & $\ci{10.2}{21.8}$ & $\ci{-1.6}{5.9}$ & 14B includes zero \\
Attenuation CI (\pp) & --- & --- & $\ci{8.1}{21.6}$ excl.\ zero \\
\bottomrule
\end{tabular}
\end{table}

The mechanism is clear: at 14B, BCP itself is ``selectively ignored''---the model maintains default behavior regardless of calibration prompts.
When BCP has no behavioral effect, DPO has nothing to interact with adversarially.
The 14B DPO model simply continues to ignore BCP, producing near-zero politeness gaps that are directionally ambiguous.

\paragraph{Other failure modes at 14B.}
\fm{1} and \fm{3} remain absent under 14B DPO, matching the base 14B BCP results: DPO does not reintroduce completion collapse or overcorrection rebound at the larger scale.
\fm{2} shows continued weak interference ($-3.8\pp$ politeness shift when targeting question rate), confirming it is the only mechanism that persists---attenuated but nonzero---across both scales and both calibration paradigms.

% ============================================================
\subsection{Full Scale-Dependent Failure Severity Matrix}
\label{sec:severity_matrix}

Table~\ref{tab:severity_matrix} synthesizes the complete taxonomy.
All four failure mechanisms exhibit capacity-dependent severity, but with different patterns: \fm{1} and \fm{3} are fully capacity-gated (catastrophic at 7B, absent at 14B), while \fm{2} and \fm{4} are attenuated but not eliminated.

\begin{table}[t]
\centering
\caption{Scale-dependent failure severity matrix. All four mechanisms are more severe at 7B.}
\label{tab:severity_matrix}
\small
\begin{tabular}{p{3.2cm}p{3.5cm}p{3.5cm}p{3.5cm}}
\toprule
\textbf{Failure Mode} & \textbf{7B Severity} & \textbf{14B Severity} & \textbf{Type} \\
\midrule
\fm{1}: Global suppression & Catastrophic (retention 19.3\%) & Absent (retention $>$97\%) & Capacity-dep.: disappears \\
\fm{2}: Cross-dim interference & Strong (pol $-25.6\%$ rel) & Weak (pol $-8.0\%$ rel) & Partially attenuated \\
\fm{3}: Overcorrection rebound & $+7.2\pp$ (4/4 runs) & Absent (0/4 runs) & Capacity-dep.: disappears \\
\fm{4}: DPO+BCP reversal & Strong (10/12, CI excl.\ 0) & Attenuated (4/12, CI incl.\ 0) & Capacity-dep.: attenuated \\
\bottomrule
\end{tabular}
\end{table}

The critical observation: despite the attenuation or disappearance of all four failure mechanisms at 14B, \textbf{H2 still fails at both scales}.
The H2 gap-reduction criterion ($\geq$30\% reduction) passes in 0/4 runs at 7B DPO+BCP and 0/4 runs at 14B DPO+BCP.
At 14B base BCP, H2 passes in only 1/4 runs (Run 2, due to an anomalously high baseline of 0.588).
The best observed gap reduction across all conditions is 27.6\% (95\% CI $\ci{18.3}{35.2}$)---still below the 30\% threshold.

This defines the \textbf{safety--efficacy gap}: scaling resolves BCP's destructive side effects (\fm{1}, \fm{3}) and attenuates its interaction pathologies (\fm{2}, \fm{4}), but does not improve calibration effectiveness.
Scaling is necessary for safe deployment of BCP (to avoid completion collapse) but insufficient for effective calibration.

% ============================================================
\section{Discussion}
\label{sec:discussion}

\paragraph{Implications for calibration research.}
Our taxonomy constrains the design space for future calibration methods in several ways.
First, prompt-based calibration faces a scale-dependent completion--calibration trade-off (\fm{1}): at smaller scales it is destructive, at larger scales it is inert.
This parallels the persona--accuracy trade-off \citep{basil2025playing_pretend, hu2026expert_persona} but adds a scale dimension.
Second, behavioral dimensions are coupled (\fm{2})---calibrating politeness without affecting question rate or completion requires methods that can disentangle dimensions, such as per-dimension reward models or constrained decoding \citep{khan2026decoding}.
Third, DPO alignment and prompt-based calibration can interact adversarially (\fm{4}), implying that calibration methods must be evaluated in combination rather than in isolation.
Fourth, the 14B disappearance of \fm{3} (overcorrection rebound) but persistence of calibration failure suggests that the problem is deeper than perturbation magnitude---gradual or iterative calibration may avoid rebound effects but likely cannot overcome the fundamental calibration ineffectiveness.

\paragraph{The ``scaling solves it'' hypothesis is misleading.}
A na\"ive reading of our scale comparison might suggest that scaling to 14B ``solves'' calibration failures, since \fm{1} and \fm{3} disappear.
Our results show this is incorrect: scaling eliminates harms but not ineffectiveness.
The failure profile transitions from ``catastrophic over-response'' at 7B (BCP overwhelms limited capacity) to ``selective ignoring'' at 14B (the model maintains default behavior regardless of calibration instructions).
Neither profile produces effective calibration.

\paragraph{Limitations.}
Our experiments use a single model family (Qwen2.5).
Cross-family generalization (e.g., Llama, Mistral) requires models with sufficient tool-calling ability in 16-tool settings; our Llama-3-8B and Mistral-7B controls achieved 5\% and 0\% completion respectively, precluding meaningful BCP evaluation.
The scale-dependent pattern of failure severity is itself a contribution, but we acknowledge it may not hold across all model families.
Our evaluation uses the $\tau$-bench retail domain; \fm{2} and \fm{3} may be task-dependent.
Sample sizes are modest ($n{=}19$--$20$ per variant for single-run experiments, $4 \times 20{=}80$ for multi-run), though bootstrap CIs confirm statistical significance for all core claims.

% ============================================================
\section{Conclusion}
\label{sec:conclusion}

We present the first taxonomy of calibration failure modes for dialect-aware LLM user simulation, identifying four mechanisms that explain why both prompt-based (BCP) and training-based (DPO) calibration strategies fail.
All four mechanisms are more severe at 7B than 14B---bootstrap confidence intervals confirm statistical significance---yet calibration remains ineffective at both scales.
The key finding is the safety--efficacy gap: scaling eliminates destructive side effects without improving calibration effectiveness.
Every calibration failure we document is more severe at 7B than 14B, yet none of them need disappear for calibration to remain ineffective.
Scaling changes how calibration fails, not whether it fails.
These findings establish that future calibration methods for LLM user simulators must go beyond prompt injection and standard preference optimization, and they must be evaluated across scales and in combination.

% ============================================================
\bibliographystyle{plainnat}
\begin{thebibliography}{10}

\bibitem[Basil et~al.(2025)]{basil2025playing_pretend}
Basil, A., et~al.
\newblock Playing pretend: Persona prompting degrades LLM accuracy.
\newblock \emph{arXiv preprint arXiv:2512.05858}, 2025.

\bibitem[Gupta et~al.(2026)]{gupta2026intra_fairness}
Gupta, R., et~al.
\newblock Intra-fairness: Bias spillover effects in targeted alignment.
\newblock \emph{arXiv preprint arXiv:2602.16438}, 2026.

\bibitem[Hu et~al.(2026)]{hu2026expert_persona}
Hu, J., et~al.
\newblock Expert persona prompting and the accuracy trade-off.
\newblock \emph{arXiv preprint arXiv:2603.18507}, 2026.

\bibitem[Khan et~al.(2026)]{khan2026decoding}
Khan, S., et~al.
\newblock Constrained decoding for controllable text generation.
\newblock \emph{arXiv preprint}, 2026.

\bibitem[Lee et~al.(2024)]{lee2024dpo_toxicity}
Lee, H., et~al.
\newblock DPO bypasses rather than removes capabilities: Fragile toxicity alignment.
\newblock In \emph{Proceedings of ICML}, 2024.

\bibitem[Rafailov et~al.(2023)]{rafailov2023dpo}
Rafailov, R., Sharma, A., Mitchell, E., Ermon, S., Manning, C.~D., and Finn, C.
\newblock Direct preference optimization: Your language model is secretly a reward model.
\newblock In \emph{Advances in NeurIPS}, 2023.

\bibitem[Seshadri et~al.(2026)]{seshadri2026lost}
Seshadri, P., et~al.
\newblock Lost in simulation: AAVE and Indian English speakers receive worst calibration.
\newblock \emph{arXiv preprint arXiv:2601.17087}, 2026.

\bibitem[Sim2Real(2026)]{sim2real2026gap}
Sim2Real Benchmark Team.
\newblock Quantifying the sim-to-real gap in $\tau$-bench user simulation.
\newblock \emph{arXiv preprint arXiv:2603.11245}, 2026.

\bibitem[Yao et~al.(2024)]{yao2024taubench}
Yao, S., et~al.
\newblock $\tau$-bench: A benchmark for tool-augmented LLM agents.
\newblock \emph{arXiv preprint arXiv:2406.12045}, 2024.

\end{thebibliography}

\end{document}
